% Author: Adam | Module: AI Job Interview Simulator | Project: SmartCareer AI
\section{Modèle d'IA et servi par Ollama}
\label{sec:modele}

\subsection{Modèle retenu}

Le module s'appuie sur \texttt{mistral:7b-instruct-v0.2-q4\_K\_M} servi par
Ollama \texttt{0.40.11} : 7\,Md de paramètres, quantification \texttt{Q4\_K\_M},
fichier de 4\,369\,387\,721 octets (4,07~Gio), famille \texttt{llama}
(relevé sur \texttt{/api/tags}). La machine de référence est un
portable équipé d'une NVIDIA GeForce RTX 3050 Ti Laptop de 4096~Mio
(chauffeur 555.99), partagée avec le Module 1.

\textbf{Pourquoi Mistral-7B.} Trois contraintes matérielles convergent.
(1)~La VRAM disponible est de 4~Gio et doit rester partagée avec Qwen3-0.6B
\,+\,LoRA du Module 1 : un 8\,Md en \texttt{q4\_K\_M} avec son cache KV ne tient
pas, un 7\,Md en tient tout juste (§\ref{sec:perf-numgpu}). (2)~La variante
\texttt{instruct} est nécessaire : le module demande une sortie JSON conforme,
pas une complétion libre. (3)~v0.2 est la version dont le gabarit
\verb|<s>[INST] <<SYS>> ... <</SYS>> instruction [/INST]| est documenté
explicitement, ce qui est indispensable puisque le module assemble lui-même ce
gabarit (§\ref{sec:modele-raw}). La qualité relative des familles sur les
critères d'entretien n'a pas été mesurée ici : \NM{} (aucun benchmark
inter-modèles n'a été lancé).

\textbf{Pourquoi \texttt{q4\_K\_M}.} C'est le compromis mémoire/qualité standard
pour un 7\,Md sur 4~Gio : \texttt{q8\_0} (\textasciitilde{}7,7~Gio) ne tient pas
en VRAM et \texttt{q2\_K} dégrade la tenue du format JSON. L'effet de la
quantification sur la qualité des scores n'a pas été mesuré : \NM{}.

\subsection{Servi par Ollama, pas par Transformers en processus}
\label{sec:modele-ollama}

Ollama tourne comme démon à part sur \texttt{localhost:11434} ; le module ne lui
parle que par HTTP (\file{llm\_adam.py}), sans importer \texttt{torch}. Ce choix
n'est pas cosmétique : le Module 1 charge déjà Torch et son LoRA dans le
processus serveur. Un second modèle dans le même processus aurait cumulé deux
allocations CUDA sous le même limiteur, sans moyen de libérer l'une sans
l'autre ; la frontière HTTP permet au contraire un
\texttt{keep\_alive} court et un endpoint de libération explicite
(§\ref{sec:perf-vram}).

\subsection{\texttt{raw=true} et assemblage manuel du gabarit}
\label{sec:modele-raw}

\begin{lstlisting}[style=adamcode, caption={Payload envoyé à \texttt{/api/generate} (\file{llm\_adam.py}, méthode \texttt{generate\_json}, lignes 145--161) ; \texttt{num\_gpu} n'est ajouté que si \texttt{ADAM\_MISTRAL\_NUM\_GPU} $\geq$ 0.}, label={lst:payload}]
payload = {
    "model": self.model,
    "prompt": prompt,
    "raw": True,
    "stream": False,
    "format": schema,
    "keep_alive": cfg.KEEP_ALIVE,
    "options": {
        "temperature": temperature,
        "num_predict": max_tokens or cfg.MAX_NEW_TOKENS,
        "num_ctx": cfg.NUM_CTX,
        "repeat_penalty": 1.1,
    },
}
if cfg.NUM_GPU >= 0:
    payload["options"]["num_gpu"] = cfg.NUM_GPU
\end{lstlisting}

Le module n'utilise \textbf{pas} \texttt{/api/chat}. Ce endpoint applique son
propre template de conversation au message reçu ; comme les prompts du module
commencent déjà par \verb|[INST] <<SYS>>|, l'usage de \texttt{/api/chat}
aurait produit un gabarit doublement emballé --- un \texttt{[INST]} à
l'intérieur du contenu d'un autre --- auquel le modèle répond par du texte
hors sujet ou par une répétition du séparateur. \texttt{raw=true} envoie la
chaîne telle quelle : la responsabilité du formatage revient entièrement au
module, ce qui est le prix à payer pour un contrôle total (et pour la
neutralisation des injections, §\ref{sec:sec-inject}).

Conséquence symétrique, documentée dans la liste des limites : \texttt{raw=true}
ne renvoie pas de liste de contexte, donc chaque appel est autonome. Questions,
évaluation et amélioration ne partagent pas de mémoire de conversation --- ce
qui correspond exactement à l'usage (§\ref{sec:archi}), mais interdit
naturellement un « relance la question 2 en plus court ».

\subsection{\texttt{format} : schéma JSON, sans garantie de grammaire}
\label{sec:modele-format}

Le schéma Pydantic de chaque étape est sérialisé en JSON Schema aplati et passé
dans \texttt{format} (les schémas complets sont reproduits en
§\ref{sec:prompts-schemas}). Ollama convertit ce schéma en contrainte de décodage
interne, \emph{mais} la documentation Ollama précise elle-même qu'il « est
important d'instruire le modèle vers du JSON dans le prompt » et qu'il est
« idéal de passer aussi le schéma JSON en chaîne dans le prompt ». Rien dans
cette documentation, ni dans nos journaux, n'établit que la grammaire soit
effectivement appliquée à Mistral en mode \texttt{raw}.

Le module ne repose donc jamais sur \texttt{format} seul : trois garde-fous lui
sont adjoints, décrits en §\ref{sec:prompts-validation}. (1)~chaque \texttt{SYSTEM}
énonce son contrat de sortie en toutes lettres (par exemple « FORMAT : renvoie
exactement \verb|{"questions": [...]}| ») et se termine par le rappel court
\texttt{FORMAT\_REMINDER} ; (2)~le schéma complet, lui, n'est recopié en texte
que dans le prompt de réparation --- là où il sert à montrer la forme fautive ;
(3)~la sortie est validée par Pydantic, réparée une fois, puis remplacée par un
repli. Mesuré sur la calibration
(§\ref{sec:tests-calibration}) : \textbf{10 réponses sur 10 valides dès le
premier essai, 0 réparation déclenchée, 0 repli}. \texttt{format} est donc
traité comme une aide statistique, pas comme une garantie.

\subsection{Paramètres de génération}
\label{sec:modele-params}

\begin{small}
\begin{table}[H]
\centering
\caption{Températures et budgets de sortie réellement utilisés.}
\label{tab:budgets}
\begin{tabular}{@{}>{\cellflush}p{2.9cm} l l
                >{\cellflush}p{5.9cm}@{}}
\toprule
\textbf{Étape} & \textbf{Temp.} & \textbf{\texttt{num\_predict}} & \textbf{Justification} \\
\midrule
Questions (\file{service.py}:175) & 0,6 & $\min(1400,\;120+170n)$ & Variété recherchée d'une séance à l'autre ; le budget suit le nombre de questions demandées (3$\to$630, 5$\to$970, 8$\to$1400 plafonné). \\
\addlinespace
Évaluation (\file{service.py}:223) & 0,2 & 420 & Stabilité : deux évaluations d'une même réponse doivent converger. Température basse, sortie bornée. \\
\addlinespace
Amélioration (\file{service.py}:271) & 0,4 & 520 & Compromis : reformulation un peu créative (STAR rédigée) mais cadrée par les lacunes diagnostiquées. \\
\addlinespace
Réparation (\file{service.py}:110) & 0,1 & identique à l'appel d'origine & Tâche de reformatage pur : la créativité ne sert à rien, elle augmente le risque de dérive sémantique. \\
\bottomrule
\end{tabular}
\end{table}
\end{small}

Les deux autres options sont globales : \texttt{num\_ctx} à 8192 tokens
(\texttt{ADAM\_MISTRAL\_NUM\_CTX}) --- au-delà, le cache KV seul dévorerait la
VRAM du Module 1 --- et \texttt{repeat\_penalty} à 1,1, retenu après l'observation
de boucles de répétition sur les réponses longues en CPU.

\subsection{Budget de contexte : la contrainte qui a dimensionné le prompt}
\label{sec:modele-budget}

8192 tokens de contexte, pour un prompt d'environ 1800 tokens
(§\ref{sec:prompts-taille}), laisse une marge de génération cohérente avec les
budgets ci-dessus. Le module ne dépasse jamais cette fenêtre, et c'est un des
motifs du raccourcissement des prompts : un prompt de 8520 caractères (état
initial des questions) consommait déjà une part dominante du contexte, ce qui ralentissait
le préfill et, en CPU, allongeait la latence perçue d'autant.

\subsection{Libération de la mémoire}
\label{sec:modele-release}

\texttt{keep\_alive} vaut \texttt{5m} par défaut
(\texttt{ADAM\_MISTRAL\_KEEP\_ALIVE}) : au-delà de cinq minutes sans appel,
Ollama décharge Mistral de lui-même. La méthode \texttt{release()}
(\file{llm\_adam.py}:215--229) force la même libération plus tôt en envoyant une
requête \texttt{keep\_alive: 0} avec un prompt vide ; elle est exposée sur
\texttt{POST /interview/release} et appelée par le navigateur à la fermeture de
l'onglet. Son échec est silencieux par conception : c'est un confort, pas une
condition de correction. La limite connue (une libération pendant une génération
en cours ne libère rien) est mesurée et discutée en §\ref{sec:perf-vram}.

\subsection{Contrôle de disponibilité}

Avant tout appel, l'UI interroge un diagnostic qui s'appuie sur
\texttt{is\_available()} (\file{llm\_adam.py}:231--245) : un \texttt{GET
/api/tags} avec timeout \texttt{(5, 10)}, puis une comparaison du nom demandé,
avec repli sur le préfixe avant \texttt{:} pour accepter
\texttt{mistral:latest} si l'étiquette exacte a changé. Trois états en
découlent, traduits en messages français (§\ref{sec:sec-erreurs}) : démon
absent, modèle absent, modèle prêt.
